\documentclass[11pt,a4paper]{article}
\usepackage[T1]{fontenc}
\usepackage[utf8]{inputenc}
\usepackage{lmodern}
\usepackage{babel}
\babelprovide[import,main]{french}
\usepackage[a4paper,margin=2.1cm]{geometry}
\usepackage{amsmath,amssymb,mathtools,graphicx,booktabs,array,xcolor}
\usepackage{microtype,hyperref,fancyhdr}
\hypersetup{colorlinks=true,linkcolor=blue!45!black,citecolor=blue!45!black,urlcolor=blue!45!black,
pdftitle={Approximation neuronale et optimisation d'une décision thermique},pdfauthor={Mhamed Halhoul}}
\pagestyle{fancy}\fancyhf{}
\fancyhead[L]{\small Approximation neuronale et décision thermique}
\fancyhead[R]{\small Mhamed Halhoul}\fancyfoot[C]{\thepage}
\setlength{\headheight}{14pt}\setlength{\parskip}{5pt}\setlength{\parindent}{0pt}
\renewcommand{\arraystretch}{1.15}
\newcommand{\Mone}{\texttt{M1\_16}}\newcommand{\Mtwo}{\texttt{M2\_16}}
\newcommand{\Mthree}{\texttt{M1\_32}}\newcommand{\That}{\widehat T_\theta}
\newcommand{\Jhat}{\widehat J_{\lambda,\theta}}
\newcommand{\fig}[3]{\begin{figure}[ht]\centering
\includegraphics[width=#2\linewidth]{figures/#1.pdf}
\caption{#3}\label{fig:#1}\end{figure}}
\begin{document}
\hypersetup{pageanchor=false}
\begin{titlepage}
\thispagestyle{empty}
\definecolor{coverink}{RGB}{20,39,62}
\definecolor{coverteal}{RGB}{0,125,133}
\begin{minipage}[c]{0.67\linewidth}
\includegraphics[width=6.5cm]{../assets/logos/uae.png}
\end{minipage}\hfill
\begin{minipage}[c]{0.20\linewidth}
\raggedleft\includegraphics[height=2.0cm]{../assets/logos/fpl.png}
\end{minipage}
\vspace{0.6cm}
\begin{center}
{\Large\bfseries\color{coverink} Université Abdelmalek Essaâdi\par}
\vspace{0.13cm}
{\large Faculté Polydisciplinaire de Larache\par}
\vspace{0.5cm}
{\small Année universitaire 2026--2027\par}
\vspace{0.18cm}
{\large Master\par}
\vspace{0.1cm}
{\large\itshape Modélisation mathématique\\
et applications et apprentissages\par}
\vspace{0.14cm}
{\small M2 en cours\par}
\vspace{1.25cm}
{\small\bfseries\color{coverteal} MINI-PROJET SCIENTIFIQUE\par}
\vspace{0.3cm}
{\fontsize{27}{32}\selectfont\bfseries\color{coverink}
Approximation neuronale\\
et optimisation d'une\\
décision thermique\par}
\vspace{0.4cm}
{\large Valeur, sensibilité et qualité de décision\\
sur un benchmark analytique\par}
\vspace{0.45cm}
{\color{coverteal}\rule{0.34\linewidth}{1.2pt}\par}
\vspace{0.45cm}
{\small Module : Algorithmes métaheuristiques\par}
\vspace{0.14cm}
{\small Réseaux de neurones pour l'apprentissage et l'optimisation\par}
\end{center}
\vfill
\begin{minipage}[t]{0.45\linewidth}
{\small\color{coverteal} Réalisé par\par}
\vspace{0.12cm}
{\large\bfseries Mhamed Halhoul\par}
\end{minipage}\hfill
\begin{minipage}[t]{0.45\linewidth}
\raggedleft
{\small\color{coverteal} Professeur encadrant :\par}
\vspace{0.12cm}
{\large\bfseries Pr. Jamal Daoudi\par}
\end{minipage}
\vspace{1.0cm}
\end{titlepage}
\setcounter{page}{1}
\hypersetup{pageanchor=true}
\subsection*{Résumé}
Un réseau précis en prédiction conduit-il à une bonne décision physique ?
Nous étudions un choix scalaire de durée de refroidissement, avec un oracle
thermique analytique fixé. Un réseau tanh $1$--$16$--$1$ approxime une température.
M1 apprend les valeurs ; M2 apprend valeurs et dérivées physiques, avec des
échelles fixes. La réplication comporte vingt initialisations appariées,
trois conditions et 3000 mises à jour Adam par trajectoire.

Le score décisionnel est la moyenne de trois différences de coût contre une
référence flottante, sans écrêtage. Les deux contrastes moyens appariés sont
favorables à M2 à 3000, avec des intervalles bootstrap entièrement négatifs.
Treize paires sur vingt sont favorables dans chaque comparaison.
Dans le panel principal M1 à 16 valeurs, les deux corrélations de rang sont
positives, mais l'IC de leur différence inclut zéro : l'hypothèse d'une
association supérieure des sensibilités reste inconclusive.

Cette étude applique des méthodes et bornes classiques. Elle ne démontre
ni nouveauté d'algorithme, ni supériorité générale, ni regret numériquement
certifié. Sa portée est un seul problème scalaire sans bruit, un réseau
et un budget fixés.

\subsection*{Abstract}
We study the association of prediction and input-sensitivity errors with
physical decision quality. A fixed scalar thermal benchmark is approximated
by a 49-parameter tanh network. Twenty paired initializations compare value-only
and value-plus-derivative supervision, including a baseline with twice as many
value labels. Mean paired decision-score differences favor derivative supervision
in this benchmark. However, the bootstrap interval for the difference between
the two primary Spearman correlations includes zero. The sensitivity-association
hypothesis remains inconclusive. Floating references and grid-error estimates
are conditional. No general superiority or algorithmic novelty is established.

\clearpage
\section{Question et portée}
\label{sec:question}
Les paramètres du réseau sont $\theta\in\mathbb R^{49}$. L'apprentissage
minimise une perte $\mathcal L(\theta)$ généralement non convexe.
Après apprentissage, la décision est une durée $p\in P=[a,b]$ pour un contexte
$\lambda$ connu :
\[
p^\star_\lambda=\arg\min_P J_\lambda(p),\qquad
\widehat p_{\lambda,\theta}\ \text{est choisi à partir de }\Jhat(p).
\]
La forte convexité physique en $p$ ne rend ni la perte convexe en $\theta$,
ni l'objectif neuronal automatiquement convexe en $p$.

Deux questions expérimentales sont séparées. La première compare l'association
de l'erreur de valeur avec la mauvaise décision à celle de l'erreur de
sensibilité. La seconde compare les stratégies de supervision, sous deux
budgets de labels distincts. Un gain moyen de M2 peut coexister avec une
différence de corrélations inconclusive.

Le regret exact est
\[
R_\lambda=J_\lambda(\widehat p)-J_\lambda(p^\star_\lambda)\geq0.
\]
Les archives utilisent une référence calculée $p_{\rm ref}$ et sauvegardent
\[
D_\lambda=J_\lambda(\widehat p)-J_\lambda(p_{\rm ref}),\qquad
S=\frac13\sum_{k\in\mathcal I}D_{\lambda_k},
\]
où $\mathcal I$ contient les trois tâches strictement intérieures fixées.
$D$ et $S$ sont des différences brutes contre une référence flottante :
ils ne sont pas le regret exact certifié. Leur signe est conservé.

La supervision de dérivées est établie, notamment dans \emph{Sobolev Training}
\cite{sobolev}. Des opérateurs informés par dérivées ont aussi des applications
à l'optimisation sous EDP \cite{dino}. L'apprentissage orienté vers une décision
\cite{spo} et la gestion d'approximations en optimisation \cite{trust} sont des
repères complémentaires. Notre réseau scalaire n'hérite pas automatiquement
des garanties ou performances de ces travaux.

L'apport réalisé est une étude expérimentale contrôlée : preuves spécialisées,
gradients analytiques, comptabilité des labels et résultats traçables.
Il n'est pas présenté comme une nouvelle méthode scientifique.

\clearpage
\section{Modèle thermique et référence M0}
\label{sec:physique}
Toutes les variables sont \textbf{sans dimension}. Une interprétation en
secondes, kelvins ou coût monétaire exigerait des échelles absentes.
La diffusivité vaut 1. Sur $0<x<1$,
\[
u(x,t)=e^{-\pi^2t}\sin(\pi x)+\frac15e^{-4\pi^2t}\sin(2\pi x)
\]
satisfait $u_t=u_{xx}$, les bords homogènes $u(0,t)=u(1,t)=0$ et
$u(x,0)=\sin(\pi x)+\frac15\sin(2\pi x)$.
Chaque dérivée temporelle et spatiale seconde multiplie le mode $k$ par
$-(k\pi)^2$. Le signe du second mode est conservé.

Au capteur $x=1/4$, pour $p\in[a,b]$, $a=0.02$, $b=0.20$ :
\[
T(p)=\frac{e^{-\pi^2p}}{\sqrt2}+\frac15e^{-4\pi^2p}.
\]
\begin{align*}
T'(p)&=-\frac{\pi^2}{\sqrt2}e^{-\pi^2p}
       -\frac{4\pi^2}{5}e^{-4\pi^2p},\\
T''(p)&=\frac{\pi^4}{\sqrt2}e^{-\pi^2p}
       +\frac{(4\pi^2)^2}{5}e^{-4\pi^2p}.
\end{align*}
On a $T'<0$, $T''>0$, $T'''<0$. Pour $J_\lambda=T+\lambda p$ :
\[
J_\lambda'=T'+\lambda,\quad J_\lambda''=T'',\quad
0<m=T''(b)\leq T''(p)\leq L_{\rm phys}=T''(a).
\]
Continuité sur le compact et forte convexité globale donnent existence et
unicité du minimiseur physique. La formule exacte démontre $m>0$,
indépendamment de $\lambda$, sans balayage de grille.
Les valeurs archivées $m\simeq9.68407$, $L_{\rm phys}\simeq198.06894$
sont des évaluations flottantes.

Posons $\lambda_b=-T'(b)$, $\lambda_a=-T'(a)$, avec $0<\lambda_b<\lambda_a$.
M0 applique
\[
p^\star_\lambda=
\begin{cases}
b,&\lambda\leq\lambda_b,\\
\text{unique racine de }T'(p)+\lambda=0,&\lambda_b<\lambda<\lambda_a,\\
a,&\lambda\geq\lambda_a.
\end{cases}
\]
Les égalités appartiennent aux frontières. À l'intérieur, une dichotomie
conserve un bracket et retourne son milieu. La plage est
$\Lambda=[\lambda_b/2,3\lambda_a/2]$.

Les références archivées utilisent une tolérance de largeur $10^{-10}$
et au plus 80 itérations. Un zéro flottant du gradient n'est pas assimilé
à une racine exacte. Les trois brackets intérieurs donnent une estimation
conditionnelle d'erreur de valeur d'environ $1.73945\,10^{-19}$ ;
ce nombre ne borne pas l'erreur flottante totale.

\clearpage
\section{Réseau et dérivée physique}
\label{sec:reseau}
Le réseau possède 16 neurones cachés tanh et une sortie linéaire :
\[
\theta=(w_0,\ldots,w_{15},\,b_0,\ldots,b_{15},\,
v_0,\ldots,v_{15},\,d).
\]
Il contient $16+16+16+1=49$ paramètres. Les normalisations fixes ne sont
pas entraînables.
\[
z=N(p)=\frac{2(p-a)}{b-a}-1=\alpha_Np+\beta_N,\quad
\alpha_N=\frac{100}{9},\quad\beta_N=-\frac{11}{9}.
\]
Le code conservé évalue $\alpha_Np+\beta_N$ en flottants, sans réajuster
les constantes ou les frontières. Pour $u_i=w_iz+b_i$, $h_i=\tanh(u_i)$,
$q_i=1-h_i^2$ :
\[
y_\theta(z)=d+\sum_i v_ih_i,\qquad
\That(p)=c_{\rm out}+s_{\rm out}y_\theta(N(p)).
\]
Les conventions sont $c_{\rm out}=0$, $s_{\rm out}=1$.
La sortie approxime directement la température sans dimension ; elle
n'est pas divisée par l'échelle de perte $S_0$.

La règle de chaîne donne
\[
\frac{dy_\theta}{dz}=\sum_i v_iw_iq_i,\qquad
\That'(p)=s_{\rm out}\alpha_N\sum_i v_iw_iq_i.
\]
Le facteur $\alpha_N$ transforme la dérivée en entrée normalisée en
dérivée par rapport à la durée physique. Le décalage $c_{\rm out}$ n'a
pas de dérivée. L'évaluation conjointe partage les seize activations.

Pour des paramètres finis, le réseau et son objectif sont continus sur $P$ :
un minimum global existe. L'architecture n'impose ni positivité,
ni monotonie, ni convexité. Une stationnarité neuronale ne démontre
pas une optimalité globale ou une bonne décision physique.

En saturation, la fonction \texttt{tanh} peut s'arrondir à $+1$ ou $-1$.
Alors $1-h_i^2$ peut valoir zéro en flottants, bien que
$\operatorname{sech}^2(u_i)>0$ pour tout $u_i$ réel fini.
Cette limite est conservée et aucun écrêtage correctif n'est ajouté.

\clearpage
\section{Pertes M1/M2 et Adam}
\label{sec:pertes}
Les échelles fixes privilégiées sont $S_0=T(a)-T(b)>0$ et $S_1=-T'(a)>0$ :
\[
S_0=0.5729505106050928,\qquad S_1=9.313693283147819.
\]
Elles proviennent des constantes archivées de l'oracle vérifié en 1B,
pas d'un ajustement sur les diagnostics.
Sur les mêmes $n$ sites pour les deux composantes de M2 :
\begin{align*}
\mathcal L_1&=\frac1{nS_0^2}\sum_j(\That(p_j)-T(p_j))^2,\\
\mathcal L_2&=\mathcal L_1+\frac1{nS_1^2}\sum_j(\That'(p_j)-T'(p_j))^2.
\end{align*}
M1 supervise les valeurs ; M2 ajoute les dérivées physiques.
Il n'y a ni facteur $1/2$, ni poids variable par point, ni régularisation.
Le coefficient de sensibilité est 1. Le $\lambda$ physique n'est pas
un coefficient de perte.

Avec $e_j=\That-T$ et $r_j=\That'-T'$ :
\[
\nabla_\theta\mathcal L_2=
\frac2{nS_0^2}\sum_j e_j\nabla_\theta\That(p_j)
+\frac2{nS_1^2}\sum_j r_j\nabla_\theta\That'(p_j).
\]
M1 conserve le premier terme. Les gradients des 49 paramètres sont
analytiques ; les expressions mixtes figurent en annexe.
La perte reçoit les labels et n'interroge pas l'oracle dans sa boucle.

Adam emploie $\eta=0.001$, $\beta_1=0.9$, $\beta_2=0.999$,
$\epsilon_A=10^{-8}$. Avec premiers moments $m_t$ et seconds moments $\nu_t$ :
\begin{align*}
m_t&=\beta_1m_{t-1}+(1-\beta_1)\nabla\mathcal L(\theta_{t-1}),\\
\nu_t&=\beta_2\nu_{t-1}+(1-\beta_2)\nabla\mathcal L(\theta_{t-1})^2,\\
\theta_t&=\theta_{t-1}
-\eta\,\frac{m_t/(1-\beta_1^t)}
{\sqrt{\nu_t/(1-\beta_2^t)}+\epsilon_A}.
\end{align*}
Les opérations sont composante par composante. $\epsilon_A$ est hors de
la racine et le compteur est global. L'état à 300 continue jusqu'à 3000
sans réinitialiser paramètres ou moments. Aucun instantané n'est choisi
sur validation et aucune convergence globale en $\theta$ n'est promise.

\clearpage
\section{Bornes classiques et hypothèses}
\label{sec:bornes}
Supposons $\sup_P|\Jhat-J_\lambda|\leq\varepsilon$ et
$\Jhat(\widehat p)\leq\inf_P\Jhat+\alpha_{\rm opt}$,
avec $\varepsilon,\alpha_{\rm opt}\geq0$. Alors
\begin{align*}
J_\lambda(\widehat p)
&\leq\Jhat(\widehat p)+\varepsilon\\
&\leq\inf_P\Jhat+\alpha_{\rm opt}+\varepsilon\\
&\leq\Jhat(p^\star)+\alpha_{\rm opt}+\varepsilon\\
&\leq J_\lambda(p^\star)+2\varepsilon+\alpha_{\rm opt}.
\end{align*}
Donc $R_\lambda\leq2\varepsilon+\alpha_{\rm opt}$.
$\alpha_{\rm opt}$ majore un défaut de minimum \emph{global} en valeur.
Une stationnarité ne le remplace pas. Une RMSE ou un maximum sur 257
sites ne remplace pas l'erreur \emph{uniforme} $\varepsilon$.

Pour $g=J_\lambda'(p)$, le résidu KKT retenu est exactement
\[
r_\lambda(p)=
\begin{cases}
|g|,&a<p<b,\\
\max(0,-g),&p=a,\\
\max(0,g),&p=b.
\end{cases}
\]
Il est la distance de zéro à $J_\lambda'(p)+N_P(p)$.
Un gradient non nul peut satisfaire l'optimalité de frontière.
Ce résidu n'est pas un résidu de gradient projeté.

La monotonie forte et la variationalité au minimiseur donnent
$m|p-p^\star|\leq r_\lambda(p)$.
La forte convexité écrite de $p$ vers $p^\star$ donne ensuite
\[
R_\lambda(p)\leq r_\lambda(p)|p-p^\star|
-\frac m2|p-p^\star|^2\leq\frac{r_\lambda(p)^2}{2m}.
\]
Ainsi
\[
|p-p^\star|\leq r_\lambda(p)/m,\qquad
R_\lambda(p)\leq r_\lambda(p)^2/(2m).
\]
Ces conclusions traitent intérieur et frontières ; elles exigent le gradient
physique et une constante $m>0$ valable sur le domaine.
Si l'erreur du gradient calculé était majorée par $\epsilon_g$, on pourrait
remplacer $r_\lambda$ par son évaluation plus $\epsilon_g$.
Sans cette majoration, un nombre flottant est un contrôle numérique,
pas un certificat. La convexité physique ne certifie pas le réseau entier.

\clearpage
\section{Référence flottante et grille globale}
\label{sec:grille}
Si le bracket M0 $[\ell,u]$ encadre la racine intérieure en arithmétique
réelle, son milieu satisfait
\[
|p_{\rm ref}-p^\star|\leq\frac{u-\ell}{2},\qquad
0\leq J_\lambda(p_{\rm ref})-J_\lambda(p^\star)
\leq\beta_{\rm ref}=\frac{L_{\rm phys}(u-\ell)^2}{8}.
\]
La seconde borne utilise la \emph{borne supérieure} de courbure et la
stationnarité intérieure. La classification de frontière donne le bon
minimiseur réel si les signes sont corrects.
Un bracket et des signes flottants ne sont pas automatiquement certifiés.

Pour $\beta_{\rm eff}=J_\lambda(p_{\rm ref})-J_\lambda(p^\star)$ :
\[
D=R-\beta_{\rm eff},\qquad R\in[D,D+\beta_{\rm ref}]
\]
sous ces hypothèses. Avec des erreurs d'évaluation majorées par $E_D$,
l'intervalle devient $[D-E_D,D+E_D+\beta_{\rm ref}]$.
Ici $E_D$ reste inconnu ; toute différence négative serait conservée.

La sélection neuronale utilise une grille de 1024 intervalles, soit 1025
candidats, avec les deux frontières. Les prédictions sont réutilisées
entre les sept $\lambda$. On choisit le meilleur objectif \emph{calculé},
avec le plus petit $p$ en cas d'égalité flottante, sans raffinement.
Température, dérivée physique et M0 ne guident pas le choix.
Ce point est un optimum de grille, pas nécessairement un point stationnaire
ou le minimum continu exact.

En arithmétique réelle,
$|\tanh''(u)|=2|h|(1-h^2)$ a pour maximum $4/(3\sqrt3)$.
Ainsi
\[
|\Jhat''|\leq\widehat B=
\frac4{3\sqrt3}|s_{\rm out}|\alpha_N^2\sum_i|v_i|w_i^2.
\]
Un minimum de frontière est candidat ; pour un minimum intérieur,
un point de grille est à distance au plus $h_{\max}/2$.
La dérivée s'y annule au minimum continu et Taylor donne
\[
\alpha_{\rm grille}\leq\widehat B h_{\max}^2/8.
\]
Aucune convexité ou unicité n'est nécessaire. Cette borne porte sur
l'objectif, pas sur la position. Son évaluation flottante n'est pas
un majorant validé par arrondis dirigés.
Si chaque objectif calculé avait une erreur majorée par $\widehat e$,
la sélection ajouterait au plus $2\widehat e$ ; cette erreur n'est pas
posée arbitrairement à zéro.

\clearpage
\section{Protocole de réplication}
\label{sec:protocole}
Le plan 2C a été figé après observation du pilote, avant la campagne 3B.
Il s'agit d'une réplication exploratoire, pas d'une préinscription antérieure
à toute donnée.
\begin{center}
\begin{tabular}{lrrrr}
\toprule
Condition & Sites & $n_T$ & $n_{T'}$ & $C_{\rm label}$ si $q_1=1$\\
\midrule
\Mone & 16 & 16 & 0 & 16\\
\Mtwo & 16 & 16 & 16 & 32\\
\Mthree & 32 & 32 & 0 & 32\\
\bottomrule
\end{tabular}
\end{center}
Les sites sont les milieux $p_j=a+(j+1/2)(b-a)/n$, pour $n=16$ ou 32.
Les labels proviennent de l'oracle analytique, évalué en flottants sans bruit.
Les 257 milieux diagnostiques sont disjoints de l'apprentissage ; les deux
frontières sont examinées séparément.

Les graines sont les entiers 20262001 à 20262020. Un générateur local tire
les 16 $w$, puis les 16 $v$, dans $[-\sqrt{6/17},\sqrt{6/17}]$ ;
les biais sont nuls. Les trois conditions reçoivent les mêmes paramètres
initiaux, avec des états Adam indépendants.
Les 60 trajectoires ont chacune 3000 mises à jour full-batch,
soit 180000 au total, et des instantanés après 300 et 3000.
Les 120 instantanés ne sont pas 120 trajectoires indépendantes.

Les sept tâches fixes sont
\[
\frac{\lambda_b}{2},\quad\lambda_b,\quad
\lambda_b+\frac{\lambda_a-\lambda_b}{4},\quad
\lambda_b+\frac{\lambda_a-\lambda_b}{2},\quad
\lambda_b+\frac{3(\lambda_a-\lambda_b)}4,\quad
\lambda_a,\quad\frac{3\lambda_a}{2}.
\]
Les trois termes centraux définissent $S$. Les quatre autres sont secondaires.
Les identités exactes du plan, pas des $\lambda$ arrondis, relient les fichiers.

\textbf{La graine est l'unité de réplication.} Chaque bloc emporte trois
conditions, deux étapes et sept tâches. Les vingt graines représentent
vingt initialisations d'un même problème physique.
Aucun arrêt sur validation, sélection de checkpoint ou réglage fondé
sur les diagnostics n'a été employé. Le résultat principal porte sur 3000 ;
l'étape 300 reste secondaire.

\clearpage
\section{Illustration fixée des réseaux sauvegardés}
\label{sec:illustration}
\fig{learned_curves}{0.98}{Valeur et sensibilité physique aux 257 sites diagnostiques.
Illustration fixée : graine 20262001, étape 3000. La référence vient du cache 3B.
Les segments relient les points, sans lissage ajouté.}
\begin{center}
\begin{tabular}{lrrr}
\toprule
Condition & RMSE $T$ & RMSE $T'$ & $S$ archivé\\
\midrule
\Mone & 0.008805 & 0.918731 & 0.004751\\
\Mtwo & 0.007184 & 0.532845 & 0.004985\\
\Mthree & 0.008726 & 0.910702 & 0.004701\\
\bottomrule
\end{tabular}
\end{center}
Pour cette graine, M2 réduit les RMSE mais son $S$ est légèrement plus élevé.
Le cas est conservé : il n'a pas été choisi pour représenter une moyenne
ou favoriser M2. Il rappelle que les métriques de prédiction et de décision
doivent être distinguées.

Cette seule illustration utilise 771 évaluations conjointes, à partir des
paramètres sauvegardés et des labels déjà acquis. Les six RMSE servant au
contrôle de lecture sont exactement égales aux archives ; la tolérance
avait été fixée à $\mathrm{abs}=10^{-12}$, $\mathrm{rel}=10^{-10}$.
Ce contrôle ne certifie pas les erreurs flottantes totales.

Aucune nouvelle acquisition physique, décision, référence M0, mise à jour
d'apprentissage ou réplication bootstrap n'est utilisée pour cette figure.
Elle est distincte des mesures de la campagne.

\clearpage
\section{Différences appariées de décision}
\label{sec:contrastes}
\fig{paired_scores}{0.88}{Vingt paires à 3000, toutes conservées.
Une liaison représente une graine. À gauche : mêmes sites et dérivées
supplémentaires ; à droite : coût de labels égal uniquement sous $q_1=1$.}
\fig{paired_effects}{0.80}{Différences moyennes de $S$ et IC bootstrap percentiles
à 95\,\%, $B=2000$, $n=20$ graines. Le signe négatif favorise M2.}
\begin{center}\small
\begin{tabular}{lrrr}
\toprule
Contraste & Moyenne & IC 95\,\% & Nég./nulles/pos.\\
\midrule
\Mtwo$-$\Mone & $-0.01006$ & $[-0.01460;-0.00557]$ & 13/0/7\\
\Mtwo$-$\Mthree & $-0.00940$ & $[-0.01390;-0.00503]$ & 13/0/7\\
\bottomrule
\end{tabular}
\end{center}
Les différences sont calculées par graine avant leur moyenne.
Les deux IC moyens sont négatifs : M2 améliore le score décisionnel moyen
dans ce benchmark. Sept différences sont positives dans chaque comparaison.
Le gain moyen ne garantit donc pas un gain pour toute graine.
Le contraste concerne la stratégie complète ; il n'identifie pas une
cause unique liée aux sensibilités.

\clearpage
\section{Association principale et réponse à l'hypothèse}
\label{sec:association}
Le panel principal est M1\_16 à 3000, vingt points sans pooling.
\[
\rho_{\rm val}=\operatorname{Spearman}(\operatorname{RMSE}_T,S),\quad
\rho_{\rm sens}=\operatorname{Spearman}(\operatorname{RMSE}_{T'},S),\quad
\Delta\rho=\rho_{\rm sens}-\rho_{\rm val}.
\]
La définition opératoire 2C emploie tout le panel et les $D$ bruts.
Elle révise explicitement la proposition historique sur le regret exact,
sans sélection a posteriori de faibles erreurs de valeur.
\fig{error_decision}{0.96}{Panel principal à 3000, $n=20$.
Même échelle verticale dans les deux panneaux. Aucune régression ou
combinaison avec les panels secondaires.}
\begin{center}
\begin{tabular}{lrr}
\toprule
Statistique & Estimation & IC bootstrap 95\,\%\\
\midrule
$\rho_{\rm val}$ & 0.9413 & $[0.8000;0.9882]$\\
$\rho_{\rm sens}$ & 0.9481 & $[0.7925;0.9958]$\\
$\Delta\rho$ & 0.0068 & $[-0.0568;0.0791]$\\
\bottomrule
\end{tabular}
\end{center}
Les deux associations sont positives. L'IC de $\Delta\rho$ inclut zéro :
\textbf{l'association supérieure de l'erreur de sensibilité reste inconclusive}.
Cette conclusion est distincte des contrastes moyens favorables à M2.

Les RMSE ont vingt valeurs distinctes chacune ; $S$ en a dix-huit,
avec deux groupes d'égalité de taille deux. Les rangs moyens utilisent
les égalités exactes des valeurs sauvegardées, sans jitter, arrondi ou
\texttt{isclose}. La robustesse numérique des classements n'est pas certifiée.

\clearpage
\section{Bootstrap et résultats secondaires}
Un générateur local de graine 20262000 construit une seule matrice de
2000 tirages de vingt indices avec remise. Chaque indice emporte tout
le bloc de graine ; les répétitions restent présentes.
La même matrice sert à tous les rho, Delta et contrastes appariés.
Les rangs sont recalculés dans chaque tirage.

Les IC sont percentiles à 95\,\%, sur les réplications définies.
Pour $m$ valeurs triées, $h=(m-1)q$ est interpolé linéairement entre
ses voisins, pour $q=0.025,0.975$.
Une variable constante donne un rho non défini, jamais zéro.
Un IC exige au moins 1800 valeurs définies sur 2000 ;
Delta exige ses deux rho. Ce seuil ne démontre pas une couverture exacte.

\begin{center}\small
\begin{tabular}{llrrrr}
\toprule
Étape & Condition & $\rho_{\rm val}$ & $\rho_{\rm sens}$ & $\Delta$ & IC de $\Delta$\\
\midrule
300 & \Mone & ND & ND & ND & ND\\
300 & \Mtwo & 0.5872 & 0.5162 & $-0.0710$ & $[-0.2792;0.0159]$\\
300 & \Mthree & ND & ND & ND & ND\\
3000 & \Mtwo & 0.9308 & 0.9714 & 0.0406 & $[-0.0045;0.1567]$\\
3000 & \Mthree & 0.9338 & 0.9383 & 0.0045 & $[-0.0484;0.0626]$\\
\bottomrule
\end{tabular}
\end{center}
À 300, $S$ est constant pour les deux conditions M1 : tous leurs tirages
rho sont non définis. Pour M2 à 300, 1976/2000 tirages sont définis, 24 ne
le sont pas. Tous les panels non constants à 3000 ont 2000/2000 tirages
définis. Aucun panel secondaire n'est promu au rang de principal.

Les deux contrastes à 300 ont la même moyenne $-0.008826$ et le même IC
$[-0.017174;-0.001916]$ : 4 différences négatives, 16 nulles.
Une moyenne appariée peut être définie quand un rho ne l'est pas.

\subsection{Frontières et transitions, hors de \texorpdfstring{$S$}{S}}
À 300, les quatre tâches secondaires ont $D=0$ pour toutes les conditions.
À 3000 :
\begin{center}\small
\begin{tabular}{lrrr}
\toprule
Tâche & \Mone & \Mtwo & \Mthree\\
\midrule
$\lambda_b/2$ : moyenne $D$ & 0.001885 & 0 & 0.001865\\
$\lambda_b/2$ : écarts de branche /20 & 6 & 0 & 6\\
$\lambda_b$ : moyenne $D$ & 0.001664 & 0.000187 & 0.001660\\
$\lambda_b$ : écarts de branche /20 & 8 & 1 & 8\\
$\lambda_a$, $3\lambda_a/2$ : moyenne $D$ & 0 & 0 & 0\\
\bottomrule
\end{tabular}
\end{center}
Les zéros de frontière ne prouvent pas une qualité intérieure.
Les intervalles secondaires restent exploratoires, sans garantie
simultanée à 95\,\%.

\clearpage
\section{Coûts et reproductibilité}
Le coût logique est $C_{\rm label}=n_T+q_1n_{T'}$.
$q_1=1$ est une convention, pas une mesure.
M2\_16 contre M1\_16 ajoute seize dérivées aux mêmes sites.
M2\_16 contre M1\_32 égalise ce coût à 32 uniquement sous cette convention,
sans égaliser sites, opérations ou durée. Les dérivées exactes de l'oracle
analytique ne deviennent pas gratuites dans une application générale.

Le cache partage les acquisitions entre graines et utilise
\texttt{quantité:p.hex()}, sans fusion approchée. Les signatures des jobs
portent sur leurs labels propres : enrichir le cache de diagnostics
ne modifie pas leur identité.
\begin{center}
\begin{tabular}{lrr}
\toprule
Compte observé 3B & $T$ & $T'$\\
\midrule
Nouveaux appels ciblés & 387 & 358\\
Demandes de cache & 1163 & 1115\\
Hits de cache & 776 & 757\\
Labels archivés distincts utilisés & 21 & 21\\
Entrées totales & 420 & 391\\
\bottomrule
\end{tabular}
\end{center}
Le cache avait importé 33 anciens labels de chaque quantité.
Les sept M0 de 2B ont été réutilisés, sans nouvel appel.
À l'import historique, deux valeurs, deux dérivées et deux courbures
ont été évaluées séparément : compte déduit du code, non directement
instrumenté. Elles ne sont pas cachées dans un total nul.

Les temps 3B observés sur ce PC sont environ 477.965 s d'apprentissage,
0.788 s de logging, 13.456 s de métriques et 43.145 s de persistance.
La durée totale de l'invocation est 566.686 s ; ces catégories ne
constituent pas une décomposition exhaustive ni un classement universel.

Les checkpoints conservent 49 paramètres, 98 moments et compteur global.
La reprise exacte a été vérifiée sur fixture avant la campagne.
Les contrôles d'intégration l'ont précédée ; oracle, réseau, gradients
et statistiques ont été vérifiés par phases. Aucune suite globale exécutée
conjointement n'est annoncée. Les 17 tests statistiques 3C ont été suivis
d'une seule analyse réelle à $B=2000$.

Les lexèmes CSV, contextes, modèles et empreintes sont archivés.
Les moyennes descriptives à 3000 sont :
\begin{center}
\begin{tabular}{lrrr}
\toprule
Condition (20 modèles) & RMSE $T$ & RMSE $T'$ & $S$\\
\midrule
\Mone & 0.011239 & 1.029386 & 0.013452\\
\Mtwo & 0.006138 & 0.442578 & 0.003389\\
\Mthree & 0.011017 & 1.015106 & 0.012788\\
\bottomrule
\end{tabular}
\end{center}
Cette présentation ne les modifie pas et ajoute uniquement l'illustration
bornée de la section~\ref{sec:illustration}.

\clearpage
\section{Discussion et conclusion}
M2 améliore le score décisionnel moyen dans les deux comparaisons prévues,
avec des IC négatifs à 3000. Le résultat concerne sa stratégie complète
dans le benchmark fixé. Il ne montre pas que réduire la RMSE de dérivée
soit l'unique cause du gain.

La question de classement reçoit une autre réponse :
$\rho_{\rm val}$ et $\rho_{\rm sens}$ sont positifs mais l'IC de leur
différence inclut zéro. L'association supérieure des sensibilités
reste inconclusive. Cela ne justifie pas un ajustement des graines
ou du budget après observation.

L'illustration fixée rappelle cette distinction : pour la graine 20262001,
M2 réduit les RMSE et augmente légèrement $S$.
Chaque comparaison moyenne comporte sept paires défavorables.
Une association in-sample n'est ni prédiction hors échantillon validée,
ni causalité, ni garantie uniforme de fiabilité.

\subsection{Limites}
Une seule variable physique, un oracle sans bruit, une architecture tanh
et un budget Adam fixés limitent la portée de l'étude.
Les vingt graines sont vingt initialisations d'un même problème,
et le protocole a été arrêté après inspection du pilote.
Les estimations de référence/grille sont conditionnelles ; les erreurs
flottantes totales ne sont pas rigoureusement majorées.
Aucun gain de vitesse sur cette formule analytique bon marché n'est revendiqué.

Les preuves spécialisent des arguments classiques.
Les références ont été lues partiellement, aux niveaux consignés dans
le registre ; aucune revue exhaustive ou nouveauté méthodologique n'en
est déduite. Les IC bootstrap empiriques, à $n=20$ et $B$ fini,
n'ont pas une couverture exacte démontrée et ne donnent pas une
probabilité bayésienne de l'hypothèse.

\subsection{Une extension future}
Une seule extension documentaire est proposée : tester le transfert de
ce protocole à un second cas physique où le coût des valeurs et des
sensibilités serait explicitement mesuré. Elle examinerait la portée
au-delà de l'oracle analytique. Modèle, budget et séparation des données
devraient être fixés avant exécution. Elle reste non exécutée.

\subsection{Conclusion}
Le pipeline relie un petit réseau à une décision physique évaluable,
avec gradients analytiques, reprise contrôlée et coûts explicites.
Dans ce benchmark, M2 améliore le score décisionnel moyen ;
l'hypothèse d'une association supérieure de l'erreur de sensibilité
reste inconclusive. L'un de ces résultats ne prouve pas l'autre.

\clearpage
\appendix
\section{Gradients mixtes et preuve contrainte}
Posons $f=\That$, $g=\That'$, $h_i=\tanh(w_iz+b_i)$, $q_i=1-h_i^2$.
$p,\alpha_N,\beta_N,c_{\rm out},s_{\rm out}$ et les cibles sont constants
lors de la dérivation en $\theta$ :
\begin{center}
\begin{tabular}{lll}
\toprule
Paramètre & $\partial f$ & $\partial g$\\
\midrule
$w_i$ & $s_{\rm out}v_i zq_i$
& $s_{\rm out}\alpha_Nv_iq_i(1-2w_i z h_i)$\\
$b_i$ & $s_{\rm out}v_iq_i$
& $-2s_{\rm out}\alpha_Nv_iw_i h_iq_i$\\
$v_i$ & $s_{\rm out}h_i$ & $s_{\rm out}\alpha_Nw_iq_i$\\
$d$ & $s_{\rm out}$ & $0$\\
\bottomrule
\end{tabular}
\end{center}
Comme $\partial_{w_i}q_i=-2h_iq_iz$ et
$\partial_{b_i}q_i=-2h_iq_i$, la dérivation de $v_iw_iq_i$
comprend le terme direct $v_iq_i$ et celui de chaîne.
$w_i=0$ n'annule donc pas généralement $\partial_{w_i}g$.
Les deux vecteurs ont 49 composantes ordonnées $w,b,v,d$.

\subsection{Complément de preuve KKT}
Le cône normal est
$N_P(p)=\{n:\langle n,y-p\rangle\leq0,\ \forall y\in P\}$.
Choisissons $n$ réalisant le résidu, $v=J'(p)+n$, $|v|=r$,
et $d=p-p^\star$.
La monotonie forte, l'optimalité au minimiseur et le signe du cône donnent
\[
m|d|^2\leq\langle J'(p)-J'(p^\star),d\rangle
\leq\langle J'(p)+n,d\rangle\leq r|d|.
\]
On utilise $\langle J'(p^\star),d\rangle\geq0$ et
$\langle n,d\rangle\geq0$. La forte convexité donne
\[
R\leq\langle J'(p),d\rangle-\frac m2|d|^2
\leq r|d|-\frac m2|d|^2\leq\frac{r^2}{2m}.
\]
La preuve respecte les frontières sans supposer un gradient nul.

\subsection{Valeurs et dérivées}
Un décalage constant $\That=T+c$ change l'erreur de valeur sans changer
les décisions. Une perturbation $\varepsilon\sin(\omega p)$ peut conserver
une petite erreur uniforme de valeur et amplifier l'erreur de dérivée.
Elle ne contredit pas $R\leq2\varepsilon+\alpha_{\rm opt}$ au minimum global.
Ces exemples classiques ne prouvent pas leur réalisation par
l'architecture fixée ni leur observation dans l'expérience.

\clearpage
\section{Sources et lecture reproductible}
Les résultats complets sont les archives 3B/3C. Les quatre données de
présentation sont dans \texttt{report/data/}. Les graphiques lisent ces CSV,
sans recalculer rho ou bootstrap. Les sources autonomes sont dans
\texttt{report/figures/}.
\begin{itemize}
\item Rapport : \texttt{report/main.tex}.
\item Diapositives : \texttt{slides/main.tex}.
\end{itemize}

La première commande d'inspection, depuis la racine, est :
\begin{verbatim}
python -I -S -B scripts/export_presentation_material.py --check
\end{verbatim}
Elle ne réévalue aucun réseau. Le contrôle documentaire est
\texttt{scripts/build\_presentation.py --check}.
Le mode de construction exige un outil local existant et la vérification
de ses options cache/hors ligne, sans apprentissage ou acquisition.

\subsection*{Niveaux de lecture}
\emph{Sobolev Training} : résumé et section 2 lus.
Le FNO informé par dérivées est une prépublication, avec lecture sélective
des sections consignées ; son acceptation n'est pas vérifiée.
Pour le cadre de région de confiance, introduction et sections 2--3 du
manuscrit auteur sont lues. Pour \emph{Smart Predict, then Optimize},
sections 2--3 de la version arXiv v5 et notice éditeur sont lues.
Aucune preuve ou implémentation de ces travaux n'est auditée intégralement.
Les preuves de ce rapport sont les arguments autonomes classiques
du dossier mathématique local.
\begin{thebibliography}{9}
\bibitem{sobolev}
Wojciech M. Czarnecki, Simon Osindero, Max Jaderberg, Grzegorz Świrszcz,
Razvan Pascanu.
\emph{Sobolev Training for Neural Networks}.
Advances in Neural Information Processing Systems, vol.~30, 2017.
\url{https://proceedings.neurips.cc/paper_files/paper/2017/hash/758a06618c69880a6cee5314ee42d52f-Abstract.html}.

\bibitem{dino}
Boyuan Yao, Dingcheng Luo, Lianghao Cao, Nikola Kovachki,
Thomas O'Leary-Roseberry, Omar Ghattas.
\emph{Derivative-Informed Fourier Neural Operator: Universal Approximation
and Applications to PDE-Constrained Optimization}.
Prépublication arXiv:2512.14086, v1 du 16 décembre 2025,
v2 du 15 mars 2026. Acceptation éditoriale non vérifiée.
\url{https://arxiv.org/abs/2512.14086}.

\bibitem{trust}
Natalia M. Alexandrov, J. E. Dennis Jr., Robert Michael Lewis, Virginia Torczon.
\emph{A trust-region framework for managing the use of approximation models
in optimization}.
Structural Optimization, vol.~15, p.~16--23, 1998.
\url{https://doi.org/10.1007/BF01197433}.

\bibitem{spo}
Adam N. Elmachtoub, Paul Grigas.
\emph{Smart \textquotedblleft Predict, then Optimize\textquotedblright}.
Management Science, vol.~68, no~1, p.~9--26, 2022
(première publication en ligne : 12 mars 2021).
\url{https://doi.org/10.1287/mnsc.2020.3922}.
\end{thebibliography}


\end{document}
